% Actual class: 02
% Slide deck: L1 Imaging v1 full
% Exact scope: pp. 25--40; L1_MATLAB.m, all sections
% NTULearn supplement: L1 Part 3 and MATLAB Demo transcripts
% Human verified: Yes

\section{第 02 节课：Extended Imaging、Perceptron 与 MATLAB Image Operations}
\label{lecture:SC4061-L02}

\lecturemeta
  {SC4061-L02}
  {2026-08-18}
  {L1 Imaging v1 full；L1\_MATLAB.m}
  {L1 pp.~25--40；MATLAB 示例全部内容}
  {课程视频 L1 Part 3 与 MATLAB Demo（字幕核对）}
  {已确认（2026-08-18）}

\subsection{本讲主线}

本讲从 traditional camera（传统相机）扩展到 computational imaging（计算成像）与其他 sensing modalities（传感模态），随后用 Perceptron（感知机）说明如何从 labelled samples（带标签样本）学习 decision boundary（决策边界）。MATLAB 部分把 image（图像）视为数组，并用 subsampling（下采样）和 logical mask（逻辑掩膜）完成最基本的图像操作。

\lecturetopic{SC4061-L02-T01}{Beyond a Conventional 2D Camera}

\begin{center}
\small
\begin{tabularx}{\textwidth}{@{}>{\raggedright\arraybackslash}p{3.5cm}>{\raggedright\arraybackslash}p{4.2cm}>{\raggedright\arraybackslash}X@{}}
  \toprule
  Imaging method（成像方式） & Measurement（测量量） & 关键结论 \\
  \midrule
  Light-field camera（光场相机） & lenslet array（微透镜阵列）采样光线的位置与方向 & 可进行 post-capture refocusing（拍摄后重聚焦），代价是数据量与空间分辨率折衷 \\
  LiDAR（激光雷达） & laser pulse（激光脉冲）的 round-trip time of flight（往返飞行时间） & 估计 depth（深度），常得到 sparse point cloud（稀疏点云） \\
  Ultrasound（超声） & 声波回波 & 不依赖可见光，可用于体内结构成像 \\
  X-ray / CT（X 射线/计算机断层成像） & X 射线穿透后的 attenuation（衰减） & radiograph（投影片）给出二维投影；CT 由多角度投影重建截面或三维体数据 \\
  MRI（磁共振成像） & magnetic resonance signal（磁共振信号） & 形成体数据，不使用 ionizing X-ray（电离 X 射线） \\
  \bottomrule
\end{tabularx}
\end{center}

LiDAR 的理想 range equation（测距公式）为
\[
  \boxed{d=\frac{c\,\Delta t}{2}},
\]
其中 $c$ 是光速，$\Delta t$ 是 round-trip time（往返时间）。除以 $2$ 是因为脉冲走过的总路程是 $2d$。

\begin{figure}[htbp]
  \centering
  \resizebox{0.82\textwidth}{!}{\begin{tikzpicture}[>=Latex,font=\small]
  \node[draw=noteBlue,rounded corners=2pt,fill=blue!5,
    minimum width=2.5cm,minimum height=1.6cm,align=center] (sensor)
    {emitter + receiver\\（发射器与接收器）};
  \node[draw=ntuRed,very thick,minimum width=0.25cm,minimum height=3.0cm,
    right=8.6cm of sensor] (target) {};
  \node[right=0.2cm of target,align=left] {object surface\\（物体表面）};

  \draw[->,very thick,draw=ntuRed]
    ([yshift=0.42cm]sensor.east) -- node[above]{emitted pulse（发射脉冲）} ([yshift=0.42cm]target.west);
  \draw[->,very thick,draw=noteBlue]
    ([yshift=-0.42cm]target.west) -- node[below]{returned pulse（返回脉冲）} ([yshift=-0.42cm]sensor.east);

  \draw[<->,draw=gray,thick]
    ([yshift=-1.45cm]sensor.east) -- node[fill=white,inner sep=2pt] {$d$} ([yshift=-1.45cm]target.west);
  \node[below=0.72cm of sensor,align=center]
    {$t_0$：send\\$t_0+\Delta t$：receive};
  \node[below=1.0cm of target,anchor=east,align=right]
    {$2d=c\Delta t$\\$d=c\Delta t/2$};
\end{tikzpicture}
}
  \caption{Time-of-Flight（飞行时间）测距：同一个传感器记录 pulse（脉冲）的发射与返回时刻。}
  \label{fig:sc4061-l02-tof}
\end{figure}

\textbf{对应例题：}\relatedexercise{SC4061-L02-E02}{LiDAR Time-of-Flight 测距}。

\lecturetopic{SC4061-L02-T02}{Perceptron as a Linear Classifier}

将 bias（偏置）吸收到增广向量中：
\[
  \tilde{\bm x}=\begin{bmatrix}\bm x\\1\end{bmatrix},
  \qquad
  \tilde{\bm w}=\begin{bmatrix}\bm w\\b\end{bmatrix},
  \qquad
  s=\tilde{\bm w}^{\mathsf T}\tilde{\bm x}.
\]
Perceptron 使用 weighted sum（加权和）与 threshold（阈值）分类：$s>0$ 判为 $C_1$，$s<0$ 判为 $C_2$；$s=0$ 构成 decision boundary：
\[
  \boxed{\tilde{\bm w}^{\mathsf T}\tilde{\bm x}=0}.
\]
二维时它是一条直线，高维时是 hyperplane（超平面）。Weights（权重）改变边界的方向，bias 改变边界的位置。

\lecturetopic{SC4061-L02-T03}{Perceptron Learning Rule}

给定 learning rate（学习率）$\alpha>0$，仅在样本被误分类或恰好落在边界上时更新：
\[
\tilde{\bm w}(k+1)=
\begin{cases}
\tilde{\bm w}(k)+\alpha\tilde{\bm x}(k),
  & \tilde{\bm x}(k)\in C_1,\ s(k)\le 0,\\
\tilde{\bm w}(k)-\alpha\tilde{\bm x}(k),
  & \tilde{\bm x}(k)\in C_2,\ s(k)\ge 0,\\
\tilde{\bm w}(k), & \text{otherwise}.
\end{cases}
\]
每次更新都会提高当前误分类样本在正确方向上的 score（得分）。但 convergence（收敛）只在 linearly-separable（线性可分）数据上有保证；即使收敛，得到的分界线也不一定具有最大 margin（间隔）。

\textbf{对应例题：}\relatedexercise{SC4061-L02-E01}{Perceptron 权重更新}。

\newpage
\lecturetopic{SC4061-L02-T04}{From Perceptron to Multi-layer Neural Networks}

Single perceptron（单个感知机）只能产生一个 linear decision boundary（线性决策边界），因此不能表示 XOR：移动或旋转一条直线都无法把对角分布的两类点完全分开。

\begin{figure}[htbp]
  \centering
  \resizebox{0.9\textwidth}{!}{\begin{tikzpicture}[>=Latex,font=\small]
  \begin{scope}
    \draw[->] (-0.25,0) -- (4.0,0) node[right] {$x_1$};
    \draw[->] (0,-0.25) -- (0,4.0) node[above] {$x_2$};
    \draw[very thick,draw=noteBlue] (0,3) -- (3,0)
      node[pos=0.58,above right,rotate=-45] {$x_1+x_2-3=0$};
    \fill[ntuRed] (3,3) circle (3pt) node[above right] {$C_1$};
    \draw[very thick,draw=noteBlue,fill=white] (1,1) circle (3pt) node[below left] {$C_2$};
    \node[align=center] at (1.9,-0.75) {linearly separable\\（线性可分）};
  \end{scope}

  \begin{scope}[xshift=7.0cm,scale=2.15]
    \draw[->] (-0.15,0) -- (1.55,0) node[right] {$x_1$};
    \draw[->] (0,-0.15) -- (0,1.55) node[above] {$x_2$};
    \fill[ntuRed] (0,1) circle (1.5pt);
    \fill[ntuRed] (1,0) circle (1.5pt);
    \draw[very thick,draw=noteBlue,fill=white] (0,0) circle (1.5pt);
    \draw[very thick,draw=noteBlue,fill=white] (1,1) circle (1.5pt);
    \draw[dashed,draw=gray] (-0.05,0.55) -- (1.2,0.55);
    \draw[dashed,draw=gray] (0.55,-0.05) -- (0.55,1.2);
    \node[scale=0.48,align=center] at (0.62,-0.38) {XOR: no single line works\\（不存在单一直线分隔）};
    \node[scale=0.48,text=ntuRed,anchor=west] at (1.18,1.18) {$C_1$: output 1};
    \node[scale=0.48,text=noteBlue,anchor=west] at (1.18,0.96) {$C_2$: output 0};
  \end{scope}
\end{tikzpicture}
}
  \caption{左：线性可分样本可由一个 Perceptron 分类；右：XOR 不存在单一直线分隔。}
  \label{fig:sc4061-l02-perceptron-xor}
\end{figure}

另一种训练思路是显式最小化 squared-error loss（平方误差损失）：
\[
  E(\bm w)=\frac12\left(r-\bm w^{\mathsf T}\bm x\right)^2,
  \qquad
  \boxed{\bm w(k+1)=\bm w(k)+\alpha
  \left[r(k)-\bm w^{\mathsf T}(k)\bm x(k)\right]\bm x(k)}.
\]
这来自 gradient descent（梯度下降）$\bm w\leftarrow\bm w-\alpha\nabla E$。与基于“是否分错”的第一种规则相比，它直接根据 prediction error（预测误差）的大小更新权重。

Multi-layer neural network（多层神经网络）把多个 neurons（神经元）分层连接，并在层间加入 nonlinear activation（非线性激活），从而组合出非线性边界。若各层之间没有非线性激活，多层线性变换仍等价于单个线性变换，不能解决 XOR。Training（训练）通常包括：
\[
\begin{aligned}
  \text{forward pass（前向传播）}
  &\longrightarrow \text{loss（损失）}\\
  &\longrightarrow \text{backpropagation（反向传播）}
  \longrightarrow \text{weight update（权重更新）}.
\end{aligned}
\]
Backpropagation 本质上用 chain rule（链式法则）高效计算各层参数的 gradient（梯度）。本课程在此只建立概念联系，不要求从零搭建完整神经网络。

\lecturetopic{SC4061-L02-T05}{MATLAB: Images Are Arrays}

MATLAB 默认以 radians（弧度）解释三角函数参数，因此
\[
  \texttt{sin(30)}\ne 0.5,
  \qquad
  \texttt{sin(30/180*pi)}=0.5.
\]
在 script（脚本）中，分号抑制命令行输出，\texttt{\%} 开始注释，\texttt{\%\%} 划分可单独运行的 section（代码节）。

\begin{center}
\small
\begin{tabularx}{\textwidth}{@{}lX@{}}
  \toprule
  MATLAB expression（表达式） & 含义 \\
  \midrule
  \texttt{a=imread('Singapore.jpg')} & 读取本地 RGB image；本次示例得到 $600\times1280\times3$ 的 \texttt{uint8} array（数组） \\
  \texttt{a(i,j,c)} & 第 $i$ 行、第 $j$ 列、第 $c$ 个 channel（通道）；索引顺序是 row, column, channel \\
  \texttt{a(i,j,:)} & 取一个 pixel（像素）的全部 RGB 分量 \\
  \texttt{a(1:s:end,1:s:end,:)} & 两个空间方向每隔 $s$ 个像素取样，保留全部通道 \\
  \bottomrule
\end{tabularx}
\end{center}

若两个方向都采用 stride（步长）$s$，高度和宽度约缩为原来的 $1/s$，pixel count（像素数）约缩为 $1/s^2$。这是 lossy downsampling（有损下采样），并不等同于一般的 image compression（图像压缩）；只有把结果重新编码保存后，文件大小才会相应变化。

\lecturetopic{SC4061-L02-T06}{Logical Mask 与 Vectorized Operation}

课程示例把满足 $R+G+B>300$ 的 bright pixels（亮像素）置为 $(60,60,60)$。先把 \texttt{uint8} image 转为 \texttt{double} 可避免整数运算的范围限制：
\[
  M(i,j)=\mathbf 1\!\left[R(i,j)+G(i,j)+B(i,j)>300\right].
\]
然后对每个 channel $c$ 统一更新：
\[
  B_c=A_c(1-M)+60M.
\]

\begin{figure}[htbp]
  \centering
  \resizebox{\textwidth}{!}{\begin{tikzpicture}[
  >=Latex,
  stage/.style={draw=noteBlue,rounded corners=2pt,fill=blue!5,
    minimum width=3.0cm,minimum height=1.35cm,align=center},
  flow/.style={->,thick,draw=noteBlue}
]
  \node[stage] (rgb) {RGB array\\$A\in\mathbb{R}^{H\times W\times3}$};
  \node[stage,right=0.65cm of rgb] (score) {brightness score\\$S=A_R+A_G+A_B$};
  \node[stage,right=0.65cm of score] (mask) {logical mask\\$M=(S>300)$};
  \node[stage,right=0.65cm of mask] (output) {output image\\$B_c=A_c(1-M)+60M$};

  \draw[flow] (rgb) -- (score);
  \draw[flow] (score) -- (mask);
  \draw[flow] (mask) -- (output);

  \node[below=0.48cm of mask,align=center]
    {$M(i,j)=1$: replace pixel\\$M(i,j)=0$: keep original};
\end{tikzpicture}
}
  \caption{Logical mask（逻辑掩膜）把 pixel-wise condition（逐像素条件）与批量赋值分开。}
  \label{fig:sc4061-l02-matlab-mask}
\end{figure}

Vectorized operation（向量化运算）把循环交给 MATLAB 的底层数组实现，通常比解释执行的双重 \texttt{for} loop 更快，也更接近公式表达；具体加速倍数取决于 MATLAB 版本、硬件和数组大小，不能把单次计时当作固定结论。

\textbf{对应例题：}\relatedexercise{SC4061-L02-E03}{RGB Threshold 与 Logical Mask}。
